\documentclass[11pt,a4paper]{article}
\usepackage[a4paper,margin=24mm]{geometry}
\usepackage[utf8]{inputenc}
\usepackage[T1]{fontenc}
\usepackage{lmodern}
\usepackage{amsmath,amssymb,amsthm}
\usepackage{booktabs,array,tabularx}
\usepackage{enumitem}
\usepackage{float}
\usepackage{longtable}
\usepackage{graphicx}
\usepackage{microtype}
\usepackage{xcolor}
\definecolor{linkblue}{rgb}{0,0.15,0.5}
\usepackage[colorlinks=true,linkcolor=linkblue,citecolor=linkblue,urlcolor=linkblue]{hyperref}
\usepackage{url}
\allowdisplaybreaks
\emergencystretch=2em

\renewcommand{\thesection}{S\arabic{section}}
\renewcommand{\theequation}{S.\arabic{equation}}
\renewcommand{\thetable}{S\arabic{section}.\arabic{table}}
\renewcommand{\thefigure}{S\arabic{section}.\arabic{figure}}
\theoremstyle{plain}
\newtheorem{theorem}{Theorem}[section]
\newtheorem{proposition}[theorem]{Proposition}
\newtheorem{lemma}[theorem]{Lemma}
\newtheorem{corollary}[theorem]{Corollary}
\theoremstyle{definition}
\newtheorem{definition}[theorem]{Definition}
\theoremstyle{remark}
\newtheorem{remark}[theorem]{Remark}

\newcommand{\R}{\mathbb{R}}
\newcommand{\Qr}{\mathbb{Q}}
\newcommand{\N}{\mathbb{N}}
\newcommand{\calA}{\mathcal{A}}
\newcommand{\calF}{\mathcal{F}}
\newcommand{\calK}{\mathcal{K}}
\newcommand{\calN}{\mathcal{N}}
\newcommand{\calP}{\mathcal{P}}
\newcommand{\calS}{\mathcal{S}}
\newcommand{\calW}{\mathcal{W}}
\newcommand{\Vstar}{V^{\ast}}
\newcommand{\Rstar}{R^{\ast}}
\newcommand{\Gstar}{G^{\ast}}
\newcommand{\tfr}[2]{\tfrac{#1}{#2}}
\newcommand{\st}[1]{\mathsf{#1}}
\newcommand{\led}{\ell}
\newcommand{\priv}{\pi}
\newcommand{\opt}{\ell^{+}}
\newcommand{\free}{\mathrm{f}}
\newcommand{\nr}{\mathrm{nr}}
\newcommand{\ofull}{o_{\mathrm{id}}}
\newcommand{\oagg}{o_{\Sigma}}
\newcommand{\oblind}{o_{\mathrm{b}}}
\newcommand{\Kinst}{K_{\mathrm{inst}}}
\newcommand{\Kagg}{K_{\mathrm{agg}}}
\newcommand{\Kcod}{K_{\mathrm{cod}}}
\newcommand{\Kfull}{K_{\mathrm{full}}}
\newcommand{\Kdec}{K_{\mathrm{dec}}}
\newcommand{\NCHECKS}{64}
\newcommand{\mref}[1]{\csname mainlab@#1\endcsname}
\newcommand{\meqref}[1]{(\mref{#1})}
\setlength{\parskip}{2pt}
\renewcommand{\arraystretch}{1.2}

\begin{document}

\title{\textbf{Supplementary material for}\\ \emph{Worked systems for viability under incomplete observation: exact belief kernels across observation structures, protocols and decentralisation, a monitoring-design census, and certified margins}}
\author{Amin Abaee\\[0.3em]{\small Independent Researcher}\\[0.2em]{\small ORCID: \href{https://orcid.org/0000-0002-0019-1842}{0000-0002-0019-1842}\quad \href{mailto:amin_abaee@ut.ac.ir}{amin\_abaee@ut.ac.ir}}}
\date{October 15, 2026}
\maketitle

\noindent This supplement collects the tables behind the counts of the main text and the material removed from it: the master table of the $36$ pairs (S1), the whole-universe census with the register conventions that are not used in the main text (S2), the stationary-policy enumeration (S3), the decentralised law tables (S4), the monitoring-design census (S5), the four-state target census (S6), the timing grids (S7), the scaling of the census to larger caps (S8), the continuous-fibre tables (S9), the verifier output and the Lean declarations (S10), and a table of symbols (S11). Theorem, section and equation numbers of the form ``Theorem~\mref{thm:mono}'' refer to the main text. Every table is generated from \texttt{results.json}.

